

\section{Introduction}

\todo[inline]{Write out better introduction, add citations. Now it is basically an extended abstract...}

Any real world measurement contains some uncertainty, and it is often crucial to understand exactly how uncertain some measurement or statement is. Large uncertainties propagate throughout theories, and hinder good predictions and foresight into the state of the system being studied. In order to make sure uncertainty is low, it must be quantifiable, measurable and comparable. Otherwise it is impossible to tell low uncertainty from large uncertainty. 

There are in essence two kinds of uncertainty in measurements: 
\begin{enumerate}
    \item the inherent uncertainty in the system, like inherent randomness;
    \item uncertainty from the lack of knowledge and information about a system. 
\end{enumerate}
The former is called \emph{aleatoric uncertainty}, or sometimes \emph{stochastic uncertainty}, while the latter -- the focus of this paper -- is called \emph{epistemic uncertainty}. 

Epistemology concerns itself with the nature of knowledge; it tries to answer questions like: what \emph{can} we know? \emph{how} can we know? and what are the limits of knowledge? This has led to a plethora of philosophical differences of opinion, as well as a wide scape of theories. In attempting to quantify these uncertainties researchers have also developed several different mathematical models, each with their pro's and con's. These usually rely on some level of semantic input, and usually describe how to combine information about these inputs. 

The first goal of this work is to define an abstract mathematical concept of \emph{epistemic syntax}. This attempts to strip away all semantic meaning, and only focus on the underlying possible mathematical logic systems for manipulating values associated to uncertainty. This puts many existing theories on a common ground, which allows us to understand their similarities and differences.  

This common language allows us to introduce mathematical descriptions of traditional epistemological philosophical positions, such as \emph{conservativity}, \emph{fallibility}, \emph{optimics}, \emph{scepticism} and \emph{idempotency}. We give these axiomatic formulations and compare the relationships between them. 

To accomplish this we use a general mathematical framework made for studying underlying and abstract structures, called \emph{category theory}. Categories are sometimes referred to as describing the mathematics \emph{behind} mathematics. It will give us an incredible flexibility, great generality, as well as opening up a wide possible world for future research. 

Having developed some understanding of the syntactical level, we use enriched category theory construct a general method for implementing an epistemic syntax into system where one wants to measure uncertainty. One upshot of this is that we naturally get an understanding of how to \emph{change} the epistemic syntax used in a system. 

\begin{introthm}
    Let $\V$ be the epistemic syntax $(0,\infty)$ with multiplication, and let $\H$ be a hypotheses category enriched in $\V$. Then the $\V$-update of $\H$ based som evidence object $E$ is Bayesian updating. 
\end{introthm}



Our approach to studying the calculus of uncertainty of course stands heavily on the shoulders of giants. Our goal is simply to contribute with a new perspective on already existing ideas, and to connect philosophical and mathematical ideas to categorical notions. 

